%%%%%%%%%%%%%%%%%%%%%%%%%%%%%%%%%%%%%%%%%%%%%%%%%%%%%%%%%%%%%%%%%%%%%%%%%%%
% Lehreinheit 1: Staunen als Erkenntnishaltung
% Quelle: WuZ-Buch, Kap. 15 (LE 1), Kap. 4 (Neurokognition),
%         Kap. 5 (Alief), Abschnitt "Ideomotor-Pendel" (10f)
% Staunmoment: Ideomotor-Pendel  |  Kernbegriff: Prediction Error
%%%%%%%%%%%%%%%%%%%%%%%%%%%%%%%%%%%%%%%%%%%%%%%%%%%%%%%%%%%%%%%%%%%%%%%%%%%
\nsl{Die erste Einheit führt nicht in eine Technik ein, sondern in eine
Haltung. Bevor ein Wort über \enquote{Wissenschaft} fällt, erleben die
Studierenden an sich selbst, dass Wahrnehmung Konstruktion ist: Ein Pendel
schwingt, und die haltende Hand hat \enquote{nichts getan}. Aus diesem
Erlebnis entwickelt die Einheit den Begriff des Prediction Errors, die
Unterscheidung von Belief und Alief und den didaktischen Dreiklang, der das
ganze Semester trägt: \emph{Staunen erzeugen -- Erklärung suchen -- Methode
verstehen}.
\newpage}
\os{\begin{description}
\item[{\rot\bf Heute}]~\\[-\lblskip]
\bi
  \ite ein Pendel, das sich \enquote{von allein} bewegt
  \ite Prediction Error: warum Überraschung lernen lässt
  \ite Belief und Alief
  \ite der Dreiklang des Semesters: Staunen -- Erklärung -- Methode
  \ite Start: Tagebuch des Staunens
\ei
\end{description}
\newpage}

\begin{vorbereitung}[Ideomotor-Pendel]
\textbf{Eine Woche vorher}
\begin{checkliste}
  \item Pendel für alle Teilnehmenden: je ein Faden von 40\,cm mit Ring,
    Mutter (M8) oder kleinem Gewicht. 40\,cm ist nach
    \textcite{cantergi2021chevreul} die Länge mit der höchsten
    Erfolgsquote. Zehn Ersatzpendel einplanen.
  \item Optional Teil~B (gemessenes Pendel): Pendel-App aus dem
    Buchprojekt (\texttt{Code/Ideomotor\_Pendel}) über HTTPS
    bereitstellen; QR-Code auf eine Folie legen. Die App braucht einmalig
    eine Sensorfreigabe auf dem Telefon, danach läuft sie offline.
  \item Selbst üben: Haltung (Ellbogen aufgestützt, Faden zwischen Daumen
    und Zeigefinger, Gewicht 1--2\,cm über der offenen Handfläche) und den
    Sprechtext so lange, bis er ohne Blick auf Notizen sitzt.
  \item Hefte für das \enquote{Tagebuch des Staunens} beschaffen (A5,
    liniert) oder die Studierenden per Mail bitten, eines mitzubringen.
\end{checkliste}

\textbf{Handout und Übungsblatt}
\begin{checkliste}
  \item Handout \enquote{Staunen als Erkenntnishaltung} und Übungsblatt
    LE~1 je einmal pro Person kopieren (Druckvorlage
    \texttt{handouts.pdf}, Seiten zu LE~1).
  \item Für die Übung \enquote{Hypothesen sortieren} je Kleingruppe ein
    A3-Blatt mit drei Spalten \emph{Staunen -- Erklärung -- Methode}
    vorbereiten.
\end{checkliste}

\textbf{Am Tag selbst}
\begin{checkliste}
  \item Tische so stellen, dass jede Person den Ellbogen aufstützen kann.
  \item Tafel in zwei Felder teilen: links \enquote{Was ist passiert?}
    (Hypothesen), rechts bleibt frei für die Auflösung.
  \item Pendel \emph{vor} der Begrüßung auf die Plätze legen -- die
    Einheit beginnt ohne Vorrede.
\end{checkliste}

\textbf{Risiken und Plan B.} Bei manchen Personen schwingt
das Pendel kaum. Das ist kein Fehlschlag, sondern ein Datum: Diese Personen
werden in der Hypothesenphase ausdrücklich gefragt, was bei ihnen anders
war. Den Effekt \emph{vor dem Debrief nicht auflösen}; die Beobachtungen mit
geschlossenen Augen und aufgelegtem Unterarm sind Daten, die zuerst
gesammelt und erst dann gedeutet werden.
\end{vorbereitung}

\begin{ablauf}
\begin{description}
\item[0--10 · Auftakt-Staunen.] Ohne Vorrede das Pendel. Achsen
  wechseln (links--rechts, vor--zurück, Kreis). Die Aha-Pause von einigen
  Sekunden bewusst aushalten.
\item[10--20 · Hypothesen sammeln.] \enquote{Was ist gerade passiert?}
  Alle Erklärungen an die Tafel, noch ohne Bewertung. Danach die beiden
  Kontrollen nach Chevreul: Augen schließen, Unterarm auflegen.
\item[20--50 · Theorie-Input.] Prediction Error und Predictive Processing;
  Belief und Alief; Staunen als epistemische Emotion; Auflösung des
  ideomotorischen Effekts und seine Aufklärungsgeschichte.
\item[50--75 · Anwendung.] Kleingruppen ordnen ihre Tafel-Hypothesen den
  Stufen \emph{Staunen -- Erklärung -- Methode} zu (Übungsblatt,
  Aufgabe~1 und 2).
\item[75--90 · Transfer und Semesterritual.] Start des Tagebuchs des
  Staunens, Auftrag zum Kurzessay, Ausblick auf den Awe Walk.
\end{description}
\end{ablauf}

%==============================================================================
\section{Staunmoment: Das Pendel}
%==============================================================================
\begin{staunbox}{Das Ideomotor-Pendel}
Halten Sie den Faden zwischen Daumen und Zeigefinger, den Ellbogen auf dem
Tisch. Das Gewicht hängt ruhig über Ihrer offenen Hand.

\emph{Bewegen Sie nichts. Denken Sie nur: links -- rechts, links -- rechts.}

\medskip
Nach 20--40 Sekunden: \emph{Jetzt vor -- zurück.} \quad
Dann: \emph{Jetzt im Kreis.}
\end{staunbox}
\os{\newpage}

\begin{description}
\item[{\rot\bf Was ist gerade passiert?}]~\\[-\lblskip]
\bi
  \ite sammeln Sie \emph{alle} Erklärungen -- ohne sie zu bewerten
  \ite Wer hat das Pendel bewegt? Woher wissen Sie das?
  \ite Bei wem hat es nicht funktioniert? Was war bei Ihnen anders?
\ei
\os{\newpage}

\item[{\rot\bf Zwei Kontrollen}]~\\[-\lblskip]
\bi
  \ite \textbf{Augen schließen} und weitermachen: Was passiert mit dem
    Ausschlag?
  \ite \textbf{Unterarm auflegen}, nur die Fingerspitzen bleiben frei:
    Was passiert jetzt?
  \ite beide Beobachtungen sind \emph{Daten} -- erst sammeln, dann deuten
\ei
\nsl{Die beiden Kontrollen gehen auf Michel-Eugène Chevreul zurück, der
das angeblich \enquote{hellsichtige} Pendel 1833 auf genau diese Weise
prüfte \parencite{chevreul1854}. Bei den meisten Personen wird der
Ausschlag mit geschlossenen Augen kleiner oder die Achse driftet, mit
aufgelegtem Unterarm bricht der Effekt weitgehend zusammen. Damit liegen
nach zehn Minuten zwei echte Befunde auf dem Tisch, die jede
Erklärung erfüllen muss.}
\os{\newpage}
\end{description}

%==============================================================================
\section{Was ist Wissenschaft?}
%==============================================================================
\nsl{Die Frage, was Wissenschaft sei, wird in dieser Einheit nicht mit einer
Definition eröffnet, sondern mit dem Erlebnis, das ihr historisch
vorausgeht. Platon lässt Sokrates im \emph{Theaitetos} sagen, das Staunen
sei der Anfang der Philosophie \parencite{platon_theaitetos}, und Aristoteles
greift den Gedanken in der \emph{Metaphysik} auf: Aus dem Staunen heraus
hätten die Menschen zu philosophieren begonnen, weil sie ihre Unwissenheit
bemerkten \parencite{aristoteles_metaphysik}.}
\begin{description}
\item[{\rot\bf Staunen als Anfang}]~\\[-\lblskip]
\bi
  \ite \emph{thaumazein}: Staunen ist der Anfang der Philosophie
    (Platon, Aristoteles)
  \ite Staunen heißt: ich \emph{bemerke}, dass mein Modell der Welt nicht
    passt
  \ite Wissenschaft ist die geordnete Form, mit diesem Bemerken umzugehen
\ei
\os{\newpage}

\item[{\rot\bf Arbeitsdefinition}]~\\
\defn{wissenschaft}{\textbf{Wissenschaft} ist ein methodisch geregeltes
Verfahren, Behauptungen über die Welt so zu prüfen, dass andere die
Prüfung nachvollziehen und wiederholen können -- und dass die Behauptung an
der Prüfung scheitern kann.}
\bi
  \ite \emph{methodisch}: das Vorgehen ist offengelegt
  \ite \emph{intersubjektiv}: andere kommen mit derselben Methode zum selben Befund
  \ite \emph{fallibel}: jedes Ergebnis ist vorläufig
\ei
\os{\newpage}
\end{description}

%==============================================================================
\section{Prediction Error: warum Überraschung lernen lässt}
%==============================================================================
\nsl{Die Neurowissenschaft beschreibt das Gehirn heute weniger als
Empfänger von Reizen denn als Vorhersagemaschine. Nach dem Ansatz des
\emph{Predictive Processing} erzeugt es fortlaufend Erwartungen darüber,
was die Sinne als Nächstes melden werden, und verarbeitet vor allem die
Abweichung zwischen Erwartung und Eingang, den \emph{Prediction Error}
\parencite{clark2016,friston2010active}. Wahrnehmung ist in diesem Bild
eine \enquote{kontrollierte Halluzination}, die von den Sinnesdaten
laufend korrigiert wird \parencite{seth2021}. Lernen findet dort statt, wo
die Korrektur nötig wird.}
\begin{description}
\item[{\rot\bf Das Gehirn als Vorhersagemaschine}]~\\[-\lblskip]
\begin{center}
\resizebox{0.85\textwidth}{!}{%
\begin{tikzpicture}[font=\sffamily,>=Stealth,
  box/.style={draw=rwuviolet,thick,fill=rwuviolet!8,rounded corners=2pt,
    minimum width=34mm,minimum height=11mm,align=center}]
  \node[box] (m) at (0,0) {Modell\\(Erwartung)};
  \node[box] (s) at (8,0) {Sinne\\(Eingang)};
  \node[box,draw=rwucyan,fill=rwucyan40!40] (e) at (4,-2.4) {Prediction Error\\= Eingang $-$ Erwartung};
  \draw[->,thick] (m) -- node[above]{Vorhersage} (s);
  \draw[->,thick] (s) |- (e);
  \draw[->,thick,rwuviolet] (e) -| node[pos=0.75,left,align=right]{Modell anpassen\\= Lernen} (m);
\end{tikzpicture}}
\end{center}
\os{\newpage}

\item[{\rot\bf Was das für das Pendel heißt}]~\\[-\lblskip]
\bi
  \ite Erwartung: \enquote{Meine Hand steht still, also bewegt sich das
    Pendel nicht.}
  \ite Eingang: das Pendel schwingt deutlich
  \ite der Fehler ist groß -- und er \emph{fühlt} sich an: das ist Staunen
  \ite Lernen beginnt erst, wenn wir das Modell ändern, statt den Befund
    wegzuerklären
\ei
\os{\newpage}

\item[{\rot\bf Staunen als epistemische Emotion}]~\\[-\lblskip]
\bi
  \ite Awe-Forschung: Staunen entsteht bei \emph{Weite} (etwas übersteigt
    den Bezugsrahmen) und \emph{Anpassungsbedarf} (das Modell muss
    geändert werden) \parencite{keltner2003}
  \ite Neugier entsteht aus einer bemerkten \emph{Informationslücke}
    \parencite{loewenstein1994}
  \ite Staunen ist damit kein Beiwerk, sondern der Moment, in dem
    Lernbereitschaft entsteht
\ei
\os{\newpage}

\item[{\rot\bf Bayesianisch gelesen}]~\\[-\lblskip]
\bi
  \ite überraschend ist, was unter den eigenen Hintergrundannahmen
    unwahrscheinlich war
  \ite eine Überraschung \emph{entwertet} diese Hintergrundannahmen --
    sie zeigt, wo das Modell revidiert werden muss
    \parencite{horwich1982probability}
  \ite Faustregel für das Semester: \textbf{Je größer die Überraschung,
    desto mehr lernt man über die eigenen Annahmen.}
\ei
\os{\newpage}
\end{description}

%==============================================================================
\section{Belief und Alief}
%==============================================================================
\nsl{Nach der Auflösung wissen alle, dass sie selbst das Pendel bewegt haben.
Trotzdem schwingt es beim nächsten Versuch wieder, und es fühlt sich wieder
so an, als täte es das von allein. Die Philosophin Tamar Gendler hat für
diesen Zustand den Begriff \emph{Alief} geprägt: eine automatische,
körpernahe Reaktionsbereitschaft, die neben der bewussten Überzeugung
(\emph{Belief}) weiterläuft und ihr widersprechen kann \parencite{gendler2008}.
Für das wissenschaftliche Arbeiten ist der Unterschied zentral: Eine
Überzeugung lässt sich durch ein Argument ändern, eine Alief nicht. Gegen
Aliefs helfen nur Verfahren, die nicht vom eigenen Erleben abhängen.}
\begin{description}
\item[{\rot\bf Zwei Arten, etwas zu \enquote{glauben}}]~\\[-\lblskip]
\begin{center}
\begin{tabular}{p{0.42\textwidth}p{0.42\textwidth}}
\toprule
\textbf{Belief} & \textbf{Alief} \\
\midrule
bewusste Überzeugung & automatische Reaktionsbereitschaft \\
\enquote{Ich weiß, dass ich das Pendel bewege.} & \enquote{Es fühlt sich an, als schwinge es von allein.} \\
durch Argumente änderbar & durch Argumente kaum änderbar \\
\bottomrule
\end{tabular}
\end{center}
\bi
  \ite Konsequenz: dem eigenen Erleben misstrauen heißt nicht, es zu
    leugnen -- sondern es zu \emph{messen}
\ei
\os{\newpage}
\end{description}

%==============================================================================
\section{Auflösung: der ideomotorische Effekt}
%==============================================================================
\nsl{Der Effekt heißt \emph{ideomotorisch}: Die bloße Vorstellung einer
Bewegung erzeugt minimale Innervationen in der Muskulatur, auch ohne
Handlungsvorsatz und unterhalb der Wahrnehmungsschwelle. Die Bewegung ist
real und messbar; nur ihre Urheberschaft wird nicht bewusst
repräsentiert. Die Aufklärungsgeschichte dieses Effekts ist zugleich eine
kleine Geschichte der experimentellen Methode.}
\begin{description}
\item[{\rot\bf Drei Schritte der Aufklärung}]~\\[-\lblskip]
\bi
  \ite \textbf{Chevreul (1833/1854):} Glasplatte, Stütze, Augenbinde --
    die Ursache liegt in der Hand, gesteuert über die Erwartung, rückgekoppelt
    über das Auge \parencite{chevreul1854}
  \ite \textbf{Carpenter (1852):} Begriff des \emph{ideomotorischen
    Prinzips} -- Vorstellungen lösen Muskelaktivität aus
    \parencite{carpenter1852}
  \ite \textbf{Faraday (1853):} ein Zeigerinstrument zum Tischrücken zeigt:
    die Hände schieben; sobald die Personen den Zeiger sehen, hört die
    Bewegung auf \parencite{faraday1853}
\ei
\os{\newpage}

\item[{\rot\bf Was wir daraus lernen}]~\\[-\lblskip]
\bi
  \ite Faraday hat nicht besser argumentiert, er hat einen \textbf{besseren
    Apparat} gebaut
  \ite Grundfigur wissenschaftlicher Kontrolle: eine Rückmeldung, die das
    Unbewusste sichtbar macht
  \ite derselbe Mechanismus steckt in Wünschelruten, Muskeltests und
    \enquote{gestützter Kommunikation} -- wir begegnen ihm in LE~4 und LE~5
    wieder
\ei
\os{\newpage}
\end{description}

%==============================================================================
\section{Der Dreiklang des Semesters}
%==============================================================================
\nsl{Jede Lehreinheit dieses Moduls folgt demselben Muster. Ein Staunmoment
erzeugt eine Informationslücke; die Gruppe sucht Erklärungen; aus der
Frage, wie man zwischen den Erklärungen entscheiden könnte, wird eine
Methode. Das Pendel hat diesen Weg heute einmal vollständig durchlaufen.}
\begin{description}
\item[{\rot\bf Staunen -- Erklärung -- Methode}]~\\[-\lblskip]
\begin{center}
\resizebox{0.9\textwidth}{!}{%
\begin{tikzpicture}[font=\sffamily,>=Stealth,
  st/.style={draw=rwuviolet,thick,fill=rwuviolet!10,rounded corners=3pt,
    minimum width=40mm,minimum height=16mm,align=center}]
  \node[st] (a) at (0,0) {\textbf{Staunen}\\Modell passt nicht};
  \node[st] (b) at (5.5,0) {\textbf{Erklärung}\\Hypothesen sammeln};
  \node[st,draw=rwucyan,fill=rwucyan40!40] (c) at (11,0) {\textbf{Methode}\\zwischen Hypothesen\\entscheiden};
  \draw[->,very thick] (a) -- (b);
  \draw[->,very thick] (b) -- (c);
  \draw[->,thick,rwugray,dashed] (c.south) to[bend left=18] node[below]{neue Befunde erzeugen neues Staunen} (a.south);
\end{tikzpicture}}
\end{center}
\bi
  \ite Pendel: \emph{Staunen} (es schwingt) -- \emph{Erklärung} (Magnetfeld?
    Luftzug? ich selbst?) -- \emph{Methode} (Augen zu, Arm auflegen, messen)
\ei
\os{\newpage}
\end{description}

\begin{merkbox}
Wissenschaft beginnt nicht mit der Antwort, sondern mit dem Bemerken, dass
das eigene Modell nicht passt. Das Gefühl, recht zu haben, ist kein Befund --
die Methode, es zu prüfen, schon.
\end{merkbox}
\os{\newpage}

%==============================================================================
\section{Übung in der Stunde}
%==============================================================================
\begin{description}
\item[{\rot\bf Hypothesen sortieren (Kleingruppen, 20 Min.)}]~\\[-\lblskip]
\bi
  \ite Übertragen Sie die Tafel-Hypothesen auf Ihr A3-Blatt
  \ite Ordnen Sie jede Aussage einer Stufe zu: Staunen, Erklärung oder
    Methode
  \ite Für jede Erklärung: Mit welcher Beobachtung könnte man sie
    \emph{widerlegen}?
\ei
\os{\newpage}
\end{description}

\begin{lehrhinweis}
Typische Hypothesen und ihre Prüfung: \enquote{Magnetismus} (Prüfung:
Pendel aus Kunststoff, Glasplatte nach Chevreul), \enquote{Luftzug}
(Prüfung: Pendel unter einer Glasglocke, Richtung wechselt mit dem
\emph{Gedanken}, nicht mit dem Raum), \enquote{ich selbst} (Prüfung:
Unterarm auflegen, App-Messung). Die Gruppe soll sehen, dass jede
Hypothese einen Test nahelegt, der sie scheitern lassen kann -- das ist die
Brücke zu LE~5.
\end{lehrhinweis}

%==============================================================================
\section{Lernziele und Aufgaben}
%==============================================================================
\begin{description}
\item[{\rot\bf Nach dieser Einheit können Sie}]~\\[-\lblskip]
\bi
  \ite den Begriff \emph{Staunen} wissenschaftstheoretisch verorten
  \ite erläutern, warum ein Prediction Error Lernen auslöst
  \ite Belief und Alief unterscheiden und an einem Beispiel erklären
  \ite eigene Vorannahmen als prüfbare Hypothesen formulieren
\ei
\os{\newpage}

\item[{\rot\bf Bis zur nächsten Woche}]~\\[-\lblskip]
\bi
  \ite \textbf{Tagebuch des Staunens} beginnen: jede Woche ein
    Staunerlebnis, verknüpft mit einem wissenschaftlichen Begriff
  \ite \textbf{Kurzessay} (500 Wörter): \enquote{Beschreiben Sie ein Erlebnis,
    das Sie zum ersten Mal an einer Überzeugung zweifeln ließ. Was hat das
    mit Wissenschaft zu tun?}
  \ite optional: \textbf{Awe Walk} -- ein Spaziergang pro Woche mit dem
    Auftrag, Dinge zu bemerken, als sähe man sie zum ersten Mal
    \parencite{sturm2022awewalk}
  \ite eine Frage mitbringen, die Sie \emph{wirklich} beschäftigt (für LE~2)
\ei
\os{\newpage}
\end{description}

\handoutinput{handouts/LE01}
